\documentclass[10pt]{article}
\usepackage[margin=20mm]{geometry}
\usepackage{amsmath,amssymb,fontspec,unicode-math,xurl,hyperref}
\IfFontExistsTF{DejaVu Serif}{\setmainfont{DejaVu Serif}}{\setmainfont{Latin Modern Roman}}
\setmathfont{latinmodern-math.otf}
\setlength{\parindent}{0pt}\setlength{\parskip}{6pt}\sloppy
\begin{document}
Abstract: The 15-page published paper, all embedded appendices and all source figures were inspected. A tenfold error in one relative improvement propagates into the headline average. The published predator-prey descriptions disagree, while correspondence between a baseline table value and its training curve is unresolved. The network specification also omits transformations needed by its stated probabilistic losses. These are Minor arithmetic and specification findings. Actual configurations and evaluation records are needed to assess the historical comparisons; no material loss of the central empirical conclusion is established. Scientific adjudication is pending.\par
Author: AI-assisted technical review\par
Date:\par
Title: Technical Verification of Ad Hoc Teamwork via Offline Goal-Based Decision Transformers\par
\# Identity and scope\par
This report reviews Xinzhi Zhang, Hohei Chan, Deheng Ye, Yi Cai and Mengchen Zhao, ICML 2025, PMLR 267, I25-43765 (0396), OpenReview `tl3FlgWScA`. All 15 physical pages were read, including Appendices A–E, Algorithms 1–3 and Tables 1–2. Figures 1–7 were actually viewed. Page references use physical numbers, matching printed pagination. The supplied manifest records no identified public implementation after bounded search, and an inaccessible separate OpenReview attachment. These are availability limits, not proof that code does not exist. No author code, pretrained teammate, dataset or simulator was run. Reviewer-written arithmetic and a small analytic diagnostic are preserved separately.\par
\# Supported findings\par
Problem 1 (The headline average includes a tenfold PP-8 error).  Table 1 (p. 13) reports TAGET 62.9 and the strongest PP-8 baseline, MADT, 62.4. The relative improvement is \[100\left(\frac{62.9}{62.4}-1\right)=0.801282\ldots\%,\] whereas the table prints \(8.01\%\). The contribution statement on p. 2 reports an average improvement of \(37.83\%\) across environments.\par
Using the strongest displayed baseline in each of the six columns gives the relative improvements \[(2.9316,\ 37.2549,\ 113.0435,\ 0.8013,\ 45.4545,\ 20.3125)\%.\] Their arithmetic mean is approximately \(36.6331\%\). Averaging the six printed percentages instead yields \(37.8317\%\), explaining the headline exactly up to rounding. This is a Minor arithmetic error: it changes the size of the gain, but leaves all six displayed point-estimate orderings unchanged. Correct the cell and headline, and state the aggregation formula and comparator explicitly. It is not appropriate to infer a different denominator from an unstated baseline merely to recover \(8.01\%\).\par
Problem 2 (The evidence does not identify one consistent experiment protocol).  Appendix A (p. 12) specifies two predator agents and four prey for PP. Figure 4 (p. 7), explicitly described as a real sampled PP trajectory, depicts one ego predator and three teammate predators. Figure 6(a) (p. 12) depicts four predators and eight numbered prey. In addition, correspondence between Figure 3’s Overcooked-8 Prompt-DT curve and Table 1 is unresolved: the table reports \(0.64\pm0.15\), whereas most of the plotted curve, including its final point, is much lower. The table does not identify a checkpoint or selection rule; an early peak is visible, so equality with the last curve point cannot be assumed.\par
These discrepancies were checked in the original source pixels, not inferred solely from extracted text. The suffixes “4” and “8” are defined in Section 5.3 as sizes of test teammate sets, so they cannot resolve the predator-count discrepancy by silently reinterpreting them as team sizes. Agent count determines the joint policy space, mirroring multiplicity, observation layout and coordination difficulty. The two-agent environment description therefore cannot unambiguously specify the experiment illustrated as an actual rollout.\par
For Overcooked-8, Table 1 makes Prompt-DT the strongest baseline, leaving a point-estimate gap of \(0.77-0.64=0.13\). Much of Figure 3 visually suggests a larger gap, but curves over training and a table of selected evaluations are not necessarily the same estimand. An old illustrative environment image, a curve from another checkpoint or a table transcription error are plausible alternatives. None is identified in the text. The figure caption and analysis still say five methods while seven are shown, which is consistent with an incompletely synchronized revision but does not establish the cause.\par
The conflicting PP descriptions establish a Minor protocol-reporting defect. The curve/table correspondence and historical quantitative impact are Needs Evidence: the source does not establish that the compared numbers use the same checkpoint, nor that a wrong team size was used in the numerical comparisons. The original Major assessment exceeded this evidence. This is not a counterexample to TAGET or an accusation about the experiments. Supply the actual PP configuration and role count, the checkpoint and evaluation identifiers for each table column and plotted curve, and the underlying per-episode returns. Then regenerate matching tables, curves and environment descriptions. If the discrepancy is confined to stale illustrations, a documented correction may suffice; if protocols differ, the affected comparisons must be rerun under a common protocol. No claim is made that correcting this issue must reverse the point-estimate ranking or invalidate every environment.\par
Problem 3 (The probabilistic output specification is incomplete).  Equations (4)–(6) (pp. 4–5) require valid Gaussian scales or covariances, and Equation (8) uses the predicted goal directly inside \(\log\widehat G\) and \(\log(1-\widehat G)\). Figure 7 (p. 15) ends the encoder and TA-Goal branches in LeakyReLU, which does not constrain variances to be positive or goal coordinates to lie in \((0,1)\). Appendix E (p. 14) also gives coefficients \(\alpha,\beta,\gamma,\sigma\) without mapping them to Equation (9)’s \(\lambda_{\rm Reg},\lambda_{\rm RTG},\lambda_{\rm TA\text{-}Goal}\) and Equation (6)’s internal weight.\par
For instance, LeakyReLU permits negative outputs and values such as 2, for which the displayed probability-space binary cross entropy is not real-valued. A sigmoid or a logits-based loss, and an exponential or softplus scale parameterization, would resolve this. They may well have been used in the actual implementation but are not specified. Thus the supported Minor finding is that the published algorithm cannot be reconstructed literally without choosing additional transformations, not that the authors’ training necessarily produced invalid numbers. Specify whether \(\sigma\) denotes standard deviation, variance or log-scale, the output transformations, the coefficient mapping, the goal feature dimensions and the treatment of \(t+k>T\). The reuse of \(i\) as both agent and goal-vector index in Equation (8) should also be clarified.\par
\# Scientific checks and bounded conclusions\par
The method’s dependency chain is explicit: global observations train the team-context encoder; stopped-gradient forward KL distils its distribution into a local proxy; the high-level modules predict return and future joint observations; a causal transformer predicts actions from local history and predicted goals. Algorithms 2–3 distinguish global training information from the proxy used at execution. In particular, the pseudocode does not show direct access to future ground-truth goals or teammates’ observations at deployment. Equation (11)’s action target is causally masked in Figure 2. No deployment information leak was established from the supplied description.\par
The local proxy cannot generally recover unobserved teammate information merely through regularization. To make this limit concrete, let two equally likely hidden team contexts produce teacher distributions \(N(-1,0.1)\) and \(N(1,0.1)\) while the ego observation is identical. The Gaussian proxy minimizing the expected forward KL is \(N(0,1.1)\), with strictly positive irreducible loss \(\tfrac12\log 11\). It represents uncertainty rather than identifying the hidden teammate. This diagnostic checks the logical reach of the mechanism; it is not a counterexample to the empirical results, because the paper assumes the observation partially reflects interactions and the low-level transformer additionally has history. Claims of completely overcoming partial observability should be interpreted within those limits. The MSE return predictor estimates a behavior-data conditional mean; no optimal-control or generalization theorem is provided.\par
Trajectory mirroring gives different ego perspectives on an existing joint trajectory. It does not create independent samples or new teammate strategies. The separated training and testing populations in Section 5.1 provide a meaningful held-out-policy test, subject to the unavailable checkpoints and configuration records. However, the described evaluation samples a teammate set for an episode; it does not specify an abrupt within-episode teammate-policy switch. Accordingly, adaptation to unfamiliar policies and changing observations is supported more directly than robust recovery from sudden policy replacement.\par
The stated confidence interval uses 50 rollout returns and a normal approximation. It measures rollout uncertainty for the evaluated model, without documenting independent training-seed variation. Under the additional assumptions that the table’s \(\pm\) values are comparable 95\% normal halfwidths and the two estimators are independent, the PP-8 difference \(0.5\) has halfwidth \(\sqrt{1.3^2+0.6^2}=1.432\), and the Overcooked-8 difference \(0.13\) has halfwidth \(\sqrt{0.15^2+0.15^2}=0.212\). Both illustrative difference intervals include zero. This is not a definitive paired significance test: pairing and covariance are unknown. It shows why point-estimate superiority alone does not establish uniform statistical significance or stability. Paired returns and training-seed replication are needed for the stronger claim.\par
Figure 5 qualitatively supports the importance of goal conditioning and mirroring, but its ablation labels do not fully define the replacement architecture when each component is removed. Table 2 supports task-dependent goal horizons (PP 6, LBF 2, Overcooked 3); it does not establish a monotonic benefit from farther prediction. The LBF best horizon value 0.153 differs from Table 1’s 0.140 and requires checkpoint/configuration identification, rather than an assumed numerical contradiction. Figure 4’s single rollout illustrates a changing predicted target; its caption’s “optimal path” is not accompanied by an optimality calculation.\par
\# Repair status\par
The arithmetic correction is independently verified. The protocol discrepancy and exact network reconstruction require author records; no invented configuration has been substituted. Assessment: Minor, with three supported local findings and the historical protocol/score impact still Needs Evidence. Fresh independent acceptance rechecked all three Problems and narrowed the original overstrong experimental-impact assessment. Adjudication remains pending and human reviews remain empty.\par
\end{document}
